% mpeg2fpga -- Guía de usuario.  Compilar con `make` (ver README.md).
\documentclass[11pt, a4paper, twoside, openany]{report}
\input{version}                 % \productversion, \productdate: generado por make
\usepackage{mpeg2fpga-style}

\begin{document}

% ======================================================================== portada
\begin{titlepage}
\begin{tikzpicture}[remember picture, overlay]
  \fill[m2navy] (current page.north west) rectangle ([yshift=-118mm]current page.north east);
  \fill[m2amber] ([yshift=-118mm]current page.north west) rectangle ([yshift=-121mm]current page.north east);
  % a macroblock grid, the decoder's unit of work
  \begin{scope}[shift={([xshift=-78mm, yshift=-18mm]current page.north east)}]
    \foreach \i in {0,...,7} \foreach \j in {0,...,5} {
      \pgfmathsetmacro{\o}{0.08 + 0.11*mod(\i*3+\j*5, 7)}
      \fill[m2teal, opacity=\o] (\i*8mm, -\j*8mm) rectangle ++(7mm, -7mm);
    }
    \draw[m2amber, line width=1.2pt] (3*8mm, -2*8mm) rectangle ++(15mm, -15mm);
  \end{scope}
  \node[anchor=north west, text=white, font=\sffamily\bfseries\fontsize{40}{44}\selectfont]
    at ([xshift=24mm, yshift=-52mm]current page.north west) {mpeg2fpga};
  \node[anchor=north west, text=white!85!m2teal, font=\sffamily\Large]
    at ([xshift=25mm, yshift=-72mm]current page.north west)
    {Decodificador MPEG-2 en PolarFire SoC};
  \node[anchor=north west, text=m2amber, font=\sffamily\large\bfseries]
    at ([xshift=25mm, yshift=-88mm]current page.north west) {Guía de usuario};
  \node[anchor=north west, font=\sffamily, text=m2grey, align=left]
    at ([xshift=25mm, yshift=-140mm]current page.north west)
    {\textbf{\color{m2navy}Versión \productversion}\\[2pt] \productdate};
  \node[anchor=south west, font=\sffamily\small, text=m2grey, align=left, text width=150mm]
    at ([xshift=25mm, yshift=24mm]current page.south west)
    {Esteban Bustamante\\ Proyecto final de Ingeniería\\[6pt]
     \scriptsize Núcleo decodificador basado en mpeg2fpga, de Koen De Vleeschauwer
     (licencia tipo BSD). Ver el apéndice de licencias.};
\end{tikzpicture}
\end{titlepage}

\pagenumbering{roman}
\tableofcontents
\clearpage
\pagenumbering{arabic}

\chapter{Introducción}

\textbf{mpeg2fpga} es un decodificador de video MPEG-2 (ISO/IEC 13818-2)
implementado en la FPGA de un PolarFire SoC y ofrecido como un servicio de red.
Un cliente le envía un clip por Ethernet y recibe, a medida que se decodifican,
los cuadros en formato YUV 4:2:0 planar (I420), en orden de presentación.

La decodificación ocurre íntegramente en hardware. Los procesadores RISC-V del
SoC sólo administran: reciben el clip, lo entregan al decodificador por DMA,
leen cada cuadro terminado y lo envían al cliente.

\section{Qué incluye un release}

\begin{dato}
\begin{tabularx}{\linewidth}{@{}l X@{}}
\textbf{Dispositivo} & Placa PolarFire SoC Discovery Kit con el bitstream del
  decodificador grabado. La FPGA es de tecnología flash: conserva la
  configuración sin alimentación ni memoria externa. \\[3pt]
\textbf{Software de la placa} & Driver de Linux (\cmd{mpeg2fpga.ko}), servicio
  \cmd{mpeg2fpgad} y su conversor nativo de cuadros. \\[3pt]
\textbf{Librería cliente} & Paquete Python \cmd{mpeg2fpga} (wheel), sin
  dependencias fuera de la biblioteca estándar. \\[3pt]
\textbf{Protocolo} & Especificación del protocolo de red v1, el contrato
  entre el dispositivo y cualquier cliente. \\[3pt]
\textbf{Documentación} & Esta guía y las notas del release con el informe de
  conformidad y rendimiento.
\end{tabularx}
\end{dato}

\section{Arquitectura}

\begin{center}
\begin{tikzpicture}[font=\sffamily\small,
    box/.style={blk, minimum width=31mm, minimum height=12mm},
    hbox/.style={hw, minimum width=31mm, minimum height=12mm}]
  % top row: software
  \node[box] (app) at (0,0) {\textbf{Aplicación}\\[-1pt]\scriptsize librería \texttt{mpeg2fpga}};
  \node[box] (d)   at (6.6,0) {\textbf{mpeg2fpgad}\\[-1pt]\scriptsize servicio, protocolo v1};
  \node[box] (drv) at (11.2,0) {\textbf{Driver}\\[-1pt]\scriptsize \texttt{mpeg2fpga.ko}};
  % bottom row: FPGA
  \node[hbox] (fs)  at (2.0,-3.4) {\textbf{Framestore}\\[-1pt]\scriptsize DDR, 4 buffers};
  \node[hbox] (dec) at (6.6,-3.4) {\textbf{Decodificador}\\[-1pt]\scriptsize MPEG-2};
  \node[hbox] (dma) at (11.2,-3.4) {\textbf{DMA de entrada}\\[-1pt]\scriptsize \texttt{stream\_dma}};
  % data path
  \draw[arr] ([yshift=2.5mm]app.east) -- node[above, lbl]{clip (HTTP)} ([yshift=2.5mm]d.west);
  \draw[arr] ([yshift=-2.5mm]d.west) -- node[below, lbl]{cuadros I420} ([yshift=-2.5mm]app.east);
  \draw[arr] (d) -- node[above, lbl]{sysfs} (drv);
  \draw[arr] (drv) -- node[right, lbl]{bloques del clip} (dma);
  \draw[arr] (dma) -- (dec);
  \draw[arr] (dec) -- node[above, lbl]{cuadros} (fs);
  % control path
  \draw[arr, dashed, color=m2teal] (dec.north east) to[bend right=12]
    node[pos=0.55, right=1mm, lbl, align=left]{IRQ:\\cuadro listo} (drv.south west);
  \draw[arr, dashed, color=m2teal] (fs.north) |- node[pos=0.3, left, lbl]{lectura} ([yshift=-6mm]d.south west) -- ([xshift=-8mm]d.south);
  \begin{scope}[on background layer]
    \node[draw=m2grey!60, dashed, rounded corners=4pt, fit=(d)(drv)(fs)(dma), inner sep=4mm] (soc) {};
    \node[lbl, anchor=south east] at (soc.north east) {PolarFire SoC};
    \node[lbl, anchor=south west] at ([yshift=6mm]app.north west) {PC del cliente};
  \end{scope}
\end{tikzpicture}
\end{center}

\begin{itemize}
  \item El \textbf{clip sube} por HTTP. El servicio lo parte en bloques y el
    driver los entrega al DMA, que alimenta al decodificador.
  \item Cuando el decodificador termina un cuadro, una \textbf{interrupción}
    le avisa al driver qué buffer del framestore lo contiene, y el driver se
    lo pasa al servicio. El servicio lo lee en el momento, antes de que el
    decodificador reutilice el buffer.
  \item El cuadro se convierte a I420 y \textbf{baja} al cliente por la misma
    conexión, mientras el resto del clip sigue subiendo.
\end{itemize}

Subida y bajada ocurren a la vez y con control de flujo en ambos sentidos: no
hay límite de largo de clip y la memoria usada en la placa es fija.

\section{Convenciones de esta guía}

Los comandos que se ejecutan en la PC del usuario no llevan prefijo; los que
se ejecutan en la placa aparecen precedidos de \cmd{ssh root@placa}. En los
ejemplos la placa tiene la dirección \cmd{192.168.18.5}.

\chapter{Especificaciones}

\section{Formatos}

\begin{center}
\begin{tabularx}{\linewidth}{@{}l X@{}}
\toprule
\textbf{Entrada} & Elementary stream de video MPEG-2 (ISO/IEC 13818-2). \\
\textbf{Perfil} & Main Profile, crominancia 4:2:0. Imágenes \emph{frame} y
  \emph{field}. \\
\textbf{Resolución} & Hasta 1920\,×\,1088 según el mapa de memoria; verificada
  en hardware hasta 720\,×\,576. \\
\textbf{Salida} & YUV 4:2:0 planar (I420), 8 bits por muestra, en orden de
  presentación. \\
\textbf{Largo de clip} & Sin límite (entrada y salida en streaming). \\
\midrule
\textbf{No soportado} & MPEG-1; escalabilidad y \emph{data partitioning};
  4:2:2 y 4:4:4; program/transport streams (el cliente debe extraer el
  elementary stream de video). \\
\bottomrule
\end{tabularx}
\end{center}

\section{Rendimiento}

Medido sobre el hardware con la herramienta de regresión del proyecto
(\cmd{tools/regress}), cinco corridas por stream luego de una de
calentamiento. El \emph{decode rate} cuenta desde el inicio del DMA hasta la
última escritura en el framestore.

\begin{center}
\begin{tabular}{@{}l c r r r@{}}
\toprule
\textbf{Stream} & \textbf{Resolución} & \textbf{Cuadros} & \textbf{fps} & \textbf{× tiempo real} \\
\midrule
tek60          & 704\,×\,480 &  60 & 23.7 & 0.79 \\
Tek-5-long     & 704\,×\,480 & 150 & 22.8 & 0.76 \\
mei-2-60f      & 704\,×\,480 &  61 & 21.1 & 0.70 \\
tcela-7        & 720\,×\,480 &  61 & 19.5 & 0.65 \\
sony-ct2       & 704\,×\,480 &  24 & 19.6 & 0.65 \\
nokia6-60      & 720\,×\,576 &  60 & 12.2 & 0.49 \\
hhi-burst-long & 720\,×\,576 &  46 & 12.9 & 0.52 \\
\bottomrule
\end{tabular}
\end{center}

\paragraph{De punta a punta} Desde la PC, con la librería cliente y el
servicio \cmd{mpeg2fpgad} sin TLS, tek60 se recibe completo a
\textbf{20\,fps} y Tek-5-long a 19.2\,fps: cerca del 85\,\% de lo que da el
decodificador solo.

\begin{nota}
El decodificador es más lento que tiempo real (0.5 a 0.8×). Para clips
grabados no es un problema: el control de flujo hace que el cliente reciba
todo, más despacio. Una fuente en vivo que no pueda frenarse, en cambio,
desbordaría cualquier buffer.
\end{nota}

\section{TLS}

TLS es opcional porque el cifrado lo hacen los procesadores de la placa, sin
extensiones criptográficas. Medido de punta a punta en la placa:

\begin{center}
\begin{tabular}{@{}l r r@{}}
\toprule
\textbf{Modo} & \textbf{MB/s} & \textbf{cuadros 704×480/s} \\
\midrule
HTTP sin cifrar            & 29.0 & $\sim$57 \\
TLS 1.2 ChaCha20-Poly1305  &  7.6 & $\sim$15 \\
TLS 1.2 AES-128-GCM        &  4.5 & $\sim$9 \\
\bottomrule
\end{tabular}
\end{center}

Con TLS la entrega queda limitada a unos 15 cuadros por segundo a 704×480. El
servicio sólo ofrece ChaCha20, el más rápido en esta CPU.

\section{Conformidad}

El decodificador se verifica contra los streams de conformidad ISO y el
decodificador de referencia (\cmd{tools/mpeg2dec}). El resultado detallado de
cada release está en sus notas; las diferencias conocidas se describen en el
capítulo~\ref{cap:limitaciones}.

\chapter{Instalación}

\section{Hardware}

\begin{itemize}
  \item Placa \textbf{PolarFire SoC Discovery Kit} (MPFS095T).
  \item Tarjeta microSD con la imagen de Linux de la placa.
  \item Cable Ethernet a la red de la PC cliente.
  \item Cable USB, sólo para programar el bitstream (JTAG) y para la consola
    serie.
\end{itemize}

Un dispositivo entregado ya trae el bitstream grabado y el software
instalado: basta conectarlo a la red. Las secciones siguientes describen la
instalación desde cero.

\section{Programar el bitstream}

El bitstream se graba por JTAG desde una PC con Microchip Libero SoC o
FlashPro Express, con el programador embebido de la placa conectado por USB.
La FPGA lo conserva sin alimentación.

\begin{lstlisting}[language=shell]
# con FlashPro Express y el archivo de programación del release
FPExpress  # abrir mpeg2fpga-<versión>.job y ejecutar PROGRAM
\end{lstlisting}

Una vez arrancado Linux, el driver informa qué bitstream está cargado:

\begin{lstlisting}[language=shell]
ssh root@placa cat /sys/bus/platform/drivers/mpeg2fpga/*/build
0.1.0+8004bc9
\end{lstlisting}

Es la versión del release y el commit del que se sintetizó; termina en
\cmd{-dirty} si se sintetizó con cambios sin registrar.

\section{Software de la placa}

Primero el driver y su overlay de device tree, después el servicio:

\begin{lstlisting}[language=shell]
scp -r driver/mpeg2fpga/tools root@placa:/tmp/m2f-driver
ssh root@placa sh /tmp/m2f-driver/install-on-board.sh

scp -r daemon api/PROTOCOL-v1.md root@placa:/tmp/m2f-daemon/
ssh root@placa sh /tmp/m2f-daemon/daemon/install-on-board.sh
\end{lstlisting}

El instalador del servicio:
\begin{itemize}
  \item crea el usuario de sistema \cmd{mpeg2fpgad} (sin privilegios);
  \item instala el servicio en \cmd{/usr/local/lib/mpeg2fpgad} y el conversor
    nativo \cmd{libm2fconv.so};
  \item genera el token de acceso en \cmd{/etc/mpeg2fpgad/token};
  \item habilita y arranca la unidad de systemd \cmd{mpeg2fpgad.service}.
\end{itemize}

Opciones: \cmd{-{}-tls} para HTTPS y \cmd{-{}-debug} para los endpoints de
diagnóstico (ver capítulo~\ref{cap:seguridad}).

\section{Verificar}

\begin{lstlisting}[language=shell]
curl http://192.168.18.5:8080/v1/health
{"ok": true}
\end{lstlisting}

\section{Librería cliente}

En la PC, con Python 3.9 o posterior:

\begin{lstlisting}[language=shell]
pip install mpeg2fpga-<versión>-py3-none-any.whl
ssh root@placa cat /etc/mpeg2fpgad/token > token
\end{lstlisting}

\chapter{Red y seguridad}\label{cap:seguridad}

La seguridad está en tres capas, cada una contra un riesgo distinto. En una
red local de confianza alcanza con la primera, que no tiene costo; las otras
permiten exponer el mismo dispositivo en una red pública sin cambiar el
protocolo.

\begin{center}
\begin{tikzpicture}[node distance=4mm]
  \node[blk, minimum width=44mm, fill=m2teal!8] (t) {\textbf{Token}\\\scriptsize ¿quién sos?};
  \node[blk, minimum width=44mm, right=of t, fill=m2blue!8] (s) {\textbf{TLS}\\\scriptsize ¿alguien escucha?};
  \node[blk, minimum width=44mm, right=of s, fill=m2amber!12] (h) {\textbf{Defensas}\\\scriptsize ¿te pueden voltear?};
  \node[lbl, below=1mm of t] {siempre · costo nulo};
  \node[lbl, below=1mm of s] {opcional · $\sim$1 core};
  \node[lbl, below=1mm of h] {siempre · costo nulo};
\end{tikzpicture}
\end{center}

\section{Token de acceso}

Todo pedido, salvo \ep{GET}{/v1/health}, debe llevar el token:

\begin{lstlisting}[language=shell]
curl -H "Authorization: Bearer $(cat token)" http://192.168.18.5:8080/v1/device
\end{lstlisting}

\begin{itemize}
  \item Es un secreto aleatorio de 256 bits, generado en la primera
    instalación, en \cmd{/etc/mpeg2fpgad/token} (sólo legible por el servicio
    y root).
  \item Se compara en tiempo constante, para no revelar información a quien
    mida cuánto tarda la respuesta.
  \item Para rotarlo:
    \cmd{ssh root@placa python3 -m mpeg2fpgad -{}-rotate-token}, y reiniciar
    el servicio.
\end{itemize}

\begin{advertencia}
Sin TLS el token viaja sin cifrar: quien pueda observar el tráfico de la red
puede copiarlo. En una red de confianza es aceptable; en una red pública, no.
\end{advertencia}

\section{TLS}

Con \cmd{-{}-tls} el servicio atiende HTTPS en el puerto 8443 con un
certificado autofirmado generado en la placa. El cliente no confía en una
autoridad certificante sino en la \emph{huella} de ese certificado, obtenida
por un canal seguro:

\begin{lstlisting}[language=shell]
ssh root@placa python3 -m mpeg2fpgad --show-fingerprint
sha256:3f9a...
\end{lstlisting}

\begin{lstlisting}[language=Python]
dev = Device("192.168.18.5", token=token, tls_fingerprint="sha256:3f9a...")
\end{lstlisting}

Si la huella no coincide, la librería se niega a conectarse
(\cmd{FingerprintMismatch}): eso es lo que impide que un tercero se haga pasar
por el dispositivo.

El costo de TLS está medido en la sección de especificaciones. Para un
despliegue público con más caudal, la alternativa es terminar TLS en un equipo
delante del dispositivo (un \emph{reverse proxy}) y dejar la placa en una red
privada.

\section{Defensas básicas}

\begin{itemize}
  \item El servicio corre como usuario sin privilegios, con el sistema de
    archivos en sólo lectura, y recibe acceso sólo a los archivos del
    decodificador que usa.
  \item Diez tokens incorrectos seguidos desde una misma dirección la bloquean
    60 segundos (respuesta 429).
  \item A lo sumo 8 conexiones simultáneas; una conexión que no envía nada en
    30 segundos se cierra.
  \item Los endpoints de diagnóstico están apagados salvo \cmd{-{}-debug}.
  \item El clip sólo llega al decodificador: nada de lo que envía el cliente
    se usa como archivo o comando. Un clip malformado puede decodificarse mal,
    pero no afecta al dispositivo.
\end{itemize}

\chapter{Uso}

\section{Librería Python}

\begin{lstlisting}[language=Python]
from mpeg2fpga import Device

dev = Device("192.168.18.5", token=open("token").read().strip())
print(dev.info().versions)          # daemon, driver, bitstream, core

result = dev.decode("clip.m2v")     # ruta, bytes o archivo abierto
with open("clip.yuv", "wb") as out:
    for frame in result:            # orden de presentación
        out.write(frame.y)
        out.write(frame.u)
        out.write(frame.v)
print(result.summary.complete, result.summary.frames_sent)
\end{lstlisting}

El archivo resultante es I420 crudo y se puede ver, por ejemplo, con
\cmd{ffplay -f rawvideo -pixel\_format yuv420p -video\_size 704x480 clip.yuv}.

\paragraph{Cada cuadro} \cmd{frame.y}, \cmd{frame.u} y \cmd{frame.v} son los
tres planos; \cmd{frame.width}, \cmd{frame.height},
\cmd{frame.display\_index}, \cmd{frame.decode\_index} y
\cmd{frame.picture\_type} (\cmd{"I"}, \cmd{"P"} o \cmd{"B"}) lo describen.
Con numpy instalado, \cmd{frame.to\_numpy()} devuelve los tres planos como
arreglos.

\paragraph{Control de flujo} El clip sube mientras se itera el resultado. Si
la aplicación se detiene entre cuadros, la subida y el decodificador también
se detienen, sin perder nada.

\paragraph{Errores}
\begin{center}
\begin{tabularx}{\linewidth}{@{}l X@{}}
\toprule
\textbf{Excepción} & \textbf{Significado} \\
\midrule
\cmd{Unauthorized} & token ausente o incorrecto \\
\cmd{DeviceBusy} & otro decode en curso; \cmd{.retry\_after} sugiere cuándo reintentar \\
\cmd{TooManyAttempts} & demasiados tokens incorrectos; esperar \\
\cmd{BadRequest} & el cuerpo no parece video MPEG-2 \\
\cmd{DecoderUnavailable} & el hardware no responde; probar \cmd{dev.reset()} \\
\cmd{DecodeIncomplete} & se cortó la conexión antes del final \\
\cmd{FingerprintMismatch} & el certificado TLS no es el esperado \\
\bottomrule
\end{tabularx}
\end{center}

\section{Desde la línea de comandos}

El protocolo es HTTP, así que alcanza con \cmd{curl}:

\begin{lstlisting}[language=shell]
curl -H "Authorization: Bearer $(cat token)" -H "Content-Type: application/octet-stream" \
     -T clip.m2v -X POST http://192.168.18.5:8080/v1/decode -o frames.bin
\end{lstlisting}

\cmd{frames.bin} contiene los registros binarios del protocolo
(capítulo~\ref{cap:protocolo}).

\section{Demo web}

\cmd{demo/demo\_server.py} es una aplicación de ejemplo construida sobre la
librería: sube un clip desde el navegador y lo reproduce a medida que llegan
los cuadros, con un resumen del decode. Corre en la PC:

\begin{lstlisting}[language=shell]
python3 demo/demo_server.py --device 192.168.18.5 --token-file token
# abrir http://localhost:8000
\end{lstlisting}

\chapter{Protocolo de red v1}\label{cap:protocolo}

Resumen del protocolo; la especificación completa está en
\cmd{PROTOCOL-v1.md}, incluida en cada release. Cualquier cliente que la
implemente funciona con el dispositivo, en cualquier lenguaje.

\section{Transporte}

HTTP/1.1 sobre TCP, puerto 8080 (8443 con TLS). Control y estado en JSON; los
cuadros en un stream binario propio. Todas las rutas llevan el prefijo
\cmd{/v1}: dentro de v1 sólo se agregan campos y endpoints, nunca se quitan.

\section{Endpoints}

\begin{center}
\begin{tabularx}{\linewidth}{@{}l X@{}}
\toprule
\ep{GET}{/v1/health} & disponible sin token; \cmd{\{"ok": true\}} \\
\ep{GET}{/v1/device} & versiones (servicio, driver, bitstream, núcleo), capacidades, límites, TLS \\
\ep{GET}{/v1/status} & estado (\cmd{idle}, \cmd{decoding}, \cmd{error}), decode en curso, última secuencia \\
\ep{POST}{/v1/decode} & el clip en el cuerpo; la respuesta es el stream de cuadros \\
\ep{POST}{/v1/reset} & recuperación: resetea el núcleo y aborta el decode en curso \\
\ep{*}{/v1/debug/...} & diagnóstico, sólo con \cmd{-{}-debug}, sin garantía de compatibilidad \\
\bottomrule
\end{tabularx}
\end{center}

Un decode a la vez: un segundo \ep{POST}{/v1/decode} recibe 409 mientras el
primero corre. Los demás endpoints responden siempre.

\section{El stream de cuadros}

Una sucesión de registros, enteros \emph{little-endian}:

\begin{center}
\begin{tikzpicture}[x=11.5mm, font=\sffamily\scriptsize,
    f/.style={draw=#1, fill=#1!10, minimum height=9mm, anchor=west, inner sep=1pt, align=center}]
  % row 1: the 12-byte header every record starts with, one unit per byte
  \foreach \x/\w/\t in {0/4/{magic\\\texttt{"M2FR"}}, 4/1/tipo, 5/1/flags,
                         6/2/{header\\\_len}, 8/4/{payload\_len}}
    \node[f=m2blue, minimum width=\w*11.5mm] at (\x, 0) {\t};
  \foreach \b in {0,4,5,6,8,12} \node[lbl, below] at (\b, -0.48) {\b};
  \draw[decorate, decoration={brace, amplitude=4pt}, m2grey]
    (0, 0.55) -- node[above=4pt, lbl]{header común: 12 bytes} (12, 0.55);
  % row 2: what follows
  \node[f=m2teal, minimum width=5*11.5mm] at (0, -1.7)
    {campos del tipo\\FRAME: 20 bytes · END: ninguno};
  \node[f=m2amber, minimum width=7*11.5mm] at (5, -1.7)
    {payload\\FRAME: cuadro I420 · END: resumen JSON};
  \draw[decorate, decoration={brace, mirror, amplitude=4pt}, m2grey]
    (0, -2.25) -- node[below=4pt, lbl]{\texttt{header\_len} $-$ 12} (5, -2.25);
  \draw[decorate, decoration={brace, mirror, amplitude=4pt}, m2grey]
    (5, -2.25) -- node[below=4pt, lbl]{\texttt{payload\_len}} (12, -2.25);
\end{tikzpicture}
\end{center}

\begin{itemize}
  \item \textbf{FRAME} (tipo 1): índice de presentación, índice de decodificación,
    ancho, alto, tipo de imagen y estructura; el payload es el cuadro en I420.
  \item \textbf{END} (tipo 2): último registro, con un resumen en JSON:
    cuadros enviados, si el decode quedó completo, bits de error del
    decodificador y tiempos por etapa.
\end{itemize}

Un cliente debe usar \cmd{header\_len} y \cmd{payload\_len} y saltear los
campos y tipos de registro que no conozca: así crece el protocolo sin romper
clientes viejos.

\section{Códigos de error}

\begin{center}
\begin{tabular}{@{}c l l@{}}
\toprule
\textbf{HTTP} & \textbf{código} & \textbf{cuándo} \\
\midrule
400 & \cmd{bad\_request} & parámetro inválido, o el cuerpo no es MPEG-2 \\
401 & \cmd{unauthorized} & token ausente o incorrecto \\
404 & \cmd{not\_found} & ruta desconocida \\
409 & \cmd{busy} & ya hay un decode en curso \\
429 & \cmd{too\_many\_attempts} & demasiados tokens incorrectos \\
503 & \cmd{decoder\_unavailable} & el hardware no responde \\
\bottomrule
\end{tabular}
\end{center}

\chapter{Operación y diagnóstico}

\section{Estado del servicio}

\begin{lstlisting}[language=shell]
ssh root@placa systemctl status mpeg2fpgad
ssh root@placa journalctl -u mpeg2fpgad -f
\end{lstlisting}

El log registra cada pedido y, por cada decode, cuántos cuadros se enviaron,
en cuánto tiempo y si quedó completo.

\section{El resumen de cada decode}

El registro END trae lo necesario para saber cómo salió un decode:

\begin{center}
\begin{tabularx}{\linewidth}{@{}l X@{}}
\toprule
\cmd{complete} & se entregó un cuadro por cada cuadro del clip \\
\cmd{decoder.error} & el decodificador encontró algo que no pudo interpretar; los cuadros pueden tener errores, el dispositivo queda sano \\
\cmd{decoder.watchdog} & el decodificador se detuvo; conviene \cmd{reset} \\
\cmd{capture\_lag\_max} & cuántos cuadros esperaron para ser leídos; 0 o 1 es normal, 2 o más indica sobrecarga \\
\cmd{timing} & tiempo por etapa en la placa: lectura, envío, espera del DMA, tiempo con el decodificador congelado \\
\bottomrule
\end{tabularx}
\end{center}

\section{Problemas frecuentes}

\begin{center}
\begin{tabularx}{\linewidth}{@{}>{\raggedright}p{42mm} X@{}}
\toprule
\textbf{Síntoma} & \textbf{Qué revisar} \\
\midrule
No responde \cmd{/v1/health} & red y dirección; \cmd{systemctl status mpeg2fpgad} \\
401 con el token correcto & el token cambió (rotación); volver a copiarlo \\
409 persistente & otro cliente está decodificando; \cmd{/v1/status} lo muestra \\
503 & driver no cargado u overlay sin aplicar: \cmd{systemctl status mpeg2fpga-overlay} \\
\cmd{complete: false} sin error & el clip termina con una imagen incompleta, o tiene imágenes que el núcleo no soporta (capítulo~\ref{cap:limitaciones}) \\
Cuadros más lentos que lo esperado & TLS activo; un cliente que lee lento (\cmd{timing.frozen\_s} alto) \\
\bottomrule
\end{tabularx}
\end{center}

\Needspace{10\baselineskip}
\section{Recuperación}

En operación normal los clips se encadenan sin resetear nada. Si un decode
termina con \cmd{watchdog}, o \cmd{/v1/status} informa \cmd{error}:

\begin{lstlisting}[language=Python]
dev.reset()
\end{lstlisting}

\chapter{Actualizaciones}

Un release tiene partes que se actualizan por caminos distintos:

\begin{center}
\begin{tabularx}{\linewidth}{@{}l l X@{}}
\toprule
\textbf{Parte} & \textbf{Dónde vive} & \textbf{Cómo se actualiza} \\
\midrule
Bitstream de la FPGA & flash de configuración de la FPGA & JTAG (FlashPro), con acceso físico \\
Driver y servicio & tarjeta microSD & por red, con los instaladores \\
Librería cliente & PC del usuario & \cmd{pip install} del nuevo wheel \\
\bottomrule
\end{tabularx}
\end{center}

\begin{advertencia}[Limitación de la placa]
El PolarFire SoC permite actualizar el bitstream en campo (\emph{Auto Update}
e \emph{IAP}), pero ambos requieren una memoria flash SPI conectada al System
Controller del chip, que la Discovery Kit no tiene. En esta placa el
bitstream sólo se actualiza por JTAG. Una placa de producto debería incluir
esa memoria.
\end{advertencia}

\section{Compatibilidad}

\begin{itemize}
  \item \cmd{/v1/device} informa las versiones de las tres partes de la placa;
    \cmd{bitstream} identifica el release y el commit exacto.
  \item Dentro del protocolo v1, un cliente de cualquier versión funciona con
    un dispositivo de cualquier versión: el protocolo sólo crece por adición.
\end{itemize}

\section{Integridad de los archivos}

Cada release publica un archivo \cmd{SHA256SUMS}. Antes de instalar:

\begin{lstlisting}[language=shell]
sha256sum -c SHA256SUMS
\end{lstlisting}

\chapter{Limitaciones conocidas}\label{cap:limitaciones}

\section{Velocidad}

El decodificador va entre 0.5 y 0.8 veces tiempo real según la resolución.
El límite está en la etapa de compensación de movimiento del núcleo, no en la
memoria ni en el reloj (medido en el proyecto). Sirve para clips grabados; no
para fuentes en vivo que no acepten control de flujo.

\section{Matrices de cuantización cargadas}

Los streams que cargan en el encabezado de secuencia una matriz de
cuantización intra no simétrica se decodifican con errores visibles desde la
primera imagen. Afecta a 11 de los 39 streams de conformidad del banco de
pruebas, entre ellos varios que cargan la matriz \emph{por defecto}. Las
matrices simétricas, y el caso habitual en que el stream no carga ninguna,
funcionan correctamente. Es una diferencia del núcleo original, en
investigación.

\section{Deriva en imágenes predichas}

En imágenes P y B con mucho movimiento la diferencia con el decodificador de
referencia crece levemente a lo largo del GOP (error medio absoluto de hasta
$\sim$1 en los streams típicos), sin efecto visible apreciable. También es
propia del núcleo original.

\section{Otras}

\begin{itemize}
  \item Sin salida de video: la placa no tiene conector de video; la salida es
    la red.
  \item MPEG-1 no está soportado; un stream \emph{stuffing} de conformidad
    (\cmd{tcela-15}) no produce imágenes.
  \item Con TLS la entrega queda limitada a unos 15 cuadros por segundo a
    704×480.
  \item Cambios de resolución dentro de un mismo clip: soportados por el
    protocolo, sin verificación exhaustiva en hardware.
\end{itemize}

\appendix
\chapter{Licencias}

\begin{center}
\begin{tabularx}{\linewidth}{@{}l l X@{}}
\toprule
\textbf{Componente} & \textbf{Licencia} & \textbf{Nota} \\
\midrule
Núcleo mpeg2fpga & tipo BSD & de Koen De Vleeschauwer; adaptado a PolarFire \\
Diseño de referencia de la placa & MIT & Microchip, base del proyecto Libero \\
Driver \cmd{mpeg2fpga.ko} & GPL-2.0 & código fuente disponible a pedido \\
Linux de la placa & GPL-2.0 & código fuente disponible a pedido \\
Servicio, librería, demo & MIT & \\
\bottomrule
\end{tabularx}
\end{center}

\section{Patentes MPEG-2}

El núcleo original incluye el aviso de licencia de patentes del pool MPEG-2
(MPEG LA). Las patentes de ese pool vencieron en la gran mayoría de los países
hacia 2018; para una comercialización conviene verificar la situación en cada
país de venta.


\end{document}
